\clearpage\thispagestyle{empty}
\pdfbookmark[0]{Quellen zur Besprechung der Connaissance des tems für 1818}{connaissance1818-sources}
\begin{center}\Large\bfseries Quellen zur Besprechung der Connaissance des tems für 1818\end{center}

Die Übertragung folgt Carl Friedrich Gauß’ Besprechung in den \emph{Werken}, Bd.~VI (Göttingen, 1874), S.~595--597. Der Erstabdruck erschien in den \emph{Göttingischen gelehrten Anzeigen}, Stück~54, 5.~April 1817, S.~531--536. Die Besprechung nennt den Jahrgang für 1818, erschienen 1815 bei der Witwe Courcier in Paris.

Im Bruch für den wahrscheinlichen Fehler von Bouvards Bestimmung der Jupitermasse lautet der Nenner im Erstabdruck (S.~535) 723. In der Werkausgabe folgt auf 72 ein undeutliches Zeichen. Die Übertragung gibt seinen Umriss wieder. Die übrigen Brüche und die Dezimalpunkte folgen der Werkausgabe. Das französische Zitat behält die Schreibweise und Hervorhebungen der Werkausgabe. Auch \emph{Connoissance} in der Besprechung von Regiomontans Aufgabe folgt der Werkausgabe.

KI-Übertragung aus dem historischen Deutsch in modernes Deutsch. KI-erstellter Satz.

\href{https://gdz.sub.uni-goettingen.de/id/PPN236007467}{Werkausgabe}. Erstdruck, S.~\href{https://images.p.sub.uni-goettingen.de/iiif/image/gdz:PPN319721507_1817:00000541/full/1650,/0/default.jpg}{531}, \href{https://images.p.sub.uni-goettingen.de/iiif/image/gdz:PPN319721507_1817:00000542/full/1650,/0/default.jpg}{532}, \href{https://images.p.sub.uni-goettingen.de/iiif/image/gdz:PPN319721507_1817:00000543/full/1650,/0/default.jpg}{533}, \href{https://images.p.sub.uni-goettingen.de/iiif/image/gdz:PPN319721507_1817:00000544/full/1650,/0/default.jpg}{534}, \href{https://images.p.sub.uni-goettingen.de/iiif/image/gdz:PPN319721507_1817:00000545/full/1650,/0/default.jpg}{535}, \href{https://images.p.sub.uni-goettingen.de/iiif/image/gdz:PPN319721507_1817:00000546/full/1650,/0/default.jpg}{536}. \href{https://cdt.imcce.fr/items/show/167}{Jahrgang für 1818, Katalog}.

\clearpage
